%pdflatex-this
\documentclass[12pt]{extarticle}[10.3.2013]
\usepackage[czech]{babel}
\usepackage[a4paper,top=1cm, bottom=1.5cm, left=1.5cm, right=1.5cm]{geometry}
\usepackage[utf8]{inputenc}
\usepackage{graphicx}
\usepackage{float}
\usepackage{multirow}
\usepackage{import}
\usepackage[unicode=true, colorlinks=true, urlcolor=twitter]{hyperref}
%\usepackage{indentfirst}
\usepackage{xcolor}
\definecolor{twitter}{HTML}{4099FF}
\definecolor{SU}{HTML}{cc3500}
\definecolor{tech}{HTML}{666666}
%\setlength{\parindent}{0in}
\usepackage[parfill]{parskip}

\author{Tomáš Kukrál}

\newcommand{\B}[1]{\textbf{{#1}}}
%\newcommand{\techSU}{\textcolor{tech}{tech@}\textcolor{SU}{SU}}
\newcommand{\techSU}{tech@SU}

\renewcommand{\familydefault}{cmss}

\begin{document}
\pagestyle{empty}

\begin{minipage}[c]{0.5\textwidth}
\begin{center}
	{\includegraphics[width=0.8\textwidth]{logo-obdelnik.png}}
\end{center}
\end{minipage}
\begin{minipage}[c]{0.5\textwidth}
\begin{center}
	{\includegraphics[width=0.8\textwidth]{logo-su.png}}
\end{center}
\end{minipage}

\part*{TECHnokrati}
\section*{Zveme Vás na povídání o studiu a životě na vysoké škole}

\B{Kdy}: 12.6.2014 od 10:00 cca do 16:00

\B{Kde}: Technická 2, Praha 6 - Dejvice, posluchárna T2:D3-309 (3. patro,
naproti schodům)



\vspace{1em}

Zveme Vás na netradiční prezentaci studia na Českém Vysokém Učení Technickém v Praze.

Na rozdíl od jiných akcí podobného typu je tato akce pořádána studenty
ve spolupráci s univerzitou, nikoliv jenom univerzitou. Díky tomu budete mít možnost zeptat se na všechny
souvislosti studia přímo studentů, kteří studiem prošli.

V dopolední části Vám nejprve představíme strukturu samotné univerzity,
ukážeme nejrůznější webové portály a důležité informační zdroje,
které vás budou provázet po celou dobu vašeho studia. 

Budete mít možnost dozvědět se informace, které se studenti obvykle
dozvídají až v průběhu studia. Díky tomu dostanete možnost se od
počátku věnovat samotnému studiu a neřešit související administrativu.


Pro odvážné jedince zde bude také příležitost naobědvat se v menze.

V odpolední části Vás stávající studenti provedou bludištěm studia na
ČVUT a upozorní Vás na nejrůznější zákoutí, která již byla mnohým osudná.

V závěru uvidíte příklady prakticky využitých znalostí, které získáte v
průběhu studia jednotlivých předmětů.

Mezi jednotlivými přednáškami vám své vysokoškolské zážitky prozradí
také úspěšní absolventi.

\vspace{3em}

\begin{minipage}[c]{0.5\textwidth}
\B{Kontakt:} \href{mailto:technokrati@su.cvut.cz}{technokrati@su.cvut.cz}, \href{https://twitter.com/klubtechsu}{@klubtechSU}, \\ \href{https://facebook.com/klubtechsu}{facebook.com/klubtechsu} \\

\B{Registrace:} \href{http://tech.su.cvut.cz/regtechnokrati}{http://tech.su.cvut.cz/regtechnokrati} \\
Registrace je důležitá, abychom věděli, kolik studentů přijede a mohli jsme rezervovat vhodné prostory. První registrovaní dostanou malou odměnu. \\

Akce je pořádána studentským klubem \techSU, který sdružuje aktivní studenty se zájmem o IT a moderní technologie. Děkujeme FEL a FIT za poskytnutí prostorů.
\end{minipage}
\begin{minipage}[c]{0.5\textwidth}
%\begin{figure}[hb]
\begin{center}
	{\includegraphics[width=0.8\textwidth]{reg-odkaz.png}}
\end{center}
%\end{figure}
\end{minipage}
\end{document}
